% Mathématiques 6e — cours V3
% Vision dans l'espace et assemblages de cubes

\section{Objectifs}

\begin{itemize}
\item Reconnaître les principaux solides usuels déjà rencontrés à l'école primaire.
\item Lire une représentation en perspective sans confondre le dessin avec le solide réel.
\item Voir dans l'espace des assemblages de cubes.
\item Passer d'un assemblage de cubes à des représentations planes : vue de dessus, de face ou de côté.
\item Reconstituer un assemblage à partir de plusieurs vues.
\item Dénombrer les cubes d'un assemblage en tenant compte des cubes cachés.
\item Décrire une stratégie de comptage claire et vérifiable.
\end{itemize}

\section{1. Solide et représentation}

Un solide est un objet géométrique de l'espace. Il possède trois dimensions : longueur, largeur et hauteur.

Une feuille ou un écran est une surface plane. Lorsqu'on y représente un solide, on produit donc une \textbf{image en deux dimensions} d'un objet qui existe en trois dimensions.

\begin{attention}
Le dessin n'est pas le solide. Une longueur dessinée plus petite sur la feuille peut représenter une arête de même longueur qu'une autre arête du solide.
\end{attention}

La représentation en perspective cherche à donner une impression de profondeur. Elle est utile pour reconnaître la structure générale d'un objet, mais elle ne conserve pas nécessairement les longueurs ni les angles du solide réel.

\section{2. Automatismes sur les solides usuels}

En 6e, on entretient la reconnaissance de solides déjà étudiés.

\begin{itemize}
\item Le \textbf{cube} possède six faces carrées.
\item Le \textbf{pavé droit} possède six faces rectangulaires.
\item Le \textbf{prisme droit} possède deux bases polygonales superposables.
\item La \textbf{pyramide} possède une base polygonale et des faces latérales triangulaires.
\item Le \textbf{cylindre} possède deux bases circulaires.
\item Le \textbf{cône} possède une base circulaire et un sommet.
\item La \textbf{boule} est un solide ; la \textbf{sphère} désigne sa surface frontière.
\end{itemize}

\begin{remarque}
Ces reconnaissances constituent surtout des automatismes. Le travail nouveau du chapitre porte sur la capacité à voir et raisonner sur des assemblages de cubes.
\end{remarque}

\section{3. Faces, arêtes et sommets}

Pour un polyèdre, on distingue trois éléments.

\begin{definition}
Une \textbf{face} est une surface plane qui limite le solide.

Une \textbf{arête} est un segment commun à deux faces.

Un \textbf{sommet} est un point où se rencontrent plusieurs arêtes.
\end{definition}

Un cube possède :
\[
6\ \text{faces},\qquad 12\ \text{arêtes},\qquad 8\ \text{sommets}.
\]

Un pavé droit possède le même nombre de faces, d'arêtes et de sommets, même si ses faces ne sont pas nécessairement carrées.

\section{4. Assemblages de cubes}

Un assemblage de cubes est obtenu en plaçant des cubes identiques les uns contre les autres, face contre face.

\coursefigure{/images/mathematiques/6eme/espace-assemblage-cubes.svg}{Assemblage de cubes unité sur plusieurs niveaux.}{Une représentation en perspective donne des informations sur les niveaux, mais certains cubes peuvent être cachés.}

Un petit cube peut servir d'\textbf{unité de comptage}. Si l'assemblage contient $N$ cubes unité, il est possible de décrire sa taille par ce nombre, même lorsque toutes les faces ne sont pas visibles.

\begin{attention}
Compter uniquement les cubes dont on voit une face entière conduit souvent à sous-estimer l'assemblage. Un cube peut être caché derrière ou sous un autre cube.
\end{attention}

\section{5. Organiser le comptage par colonnes}

Une stratégie efficace consiste à observer l'assemblage verticalement.

Supposons que les hauteurs des colonnes visibles depuis le dessus soient :
\[
2,\quad 1,\quad 3,\quad 2,\quad 1.
\]

Le nombre total de cubes vaut :
\[
2+1+3+2+1=9.
\]

\begin{methode}
Pour compter par colonnes :
\begin{enumerate}
\item repérer chaque position occupée au sol ;
\item déterminer la hauteur de la pile à cette position ;
\item additionner les hauteurs ;
\item vérifier qu'aucune position n'a été comptée deux fois.
\end{enumerate}
\end{methode}

Cette méthode est particulièrement efficace lorsqu'on dispose d'une vue de dessus accompagnée des hauteurs.

\section{6. Organiser le comptage par couches}

On peut aussi couper mentalement le solide horizontalement.

Si la couche inférieure contient $8$ cubes, la deuxième couche $5$ cubes et la troisième couche $2$ cubes, alors :
\[
N=8+5+2=15.
\]

\begin{exemple}
Un assemblage comporte une base rectangulaire complète de $4$ cubes sur $3$, puis une deuxième couche de $5$ cubes.

La première couche contient :
\[
4\times3=12
\]
cubes.

Le total vaut :
\[
12+5=17.
\]

L'assemblage contient donc $\boxed{17}$ cubes.
\end{exemple}

Le comptage par couches évite de parcourir le dessin de manière désordonnée.

\section{7. Les cubes cachés}

Dans une perspective, un cube situé derrière un autre peut ne laisser apparaître aucune face.

Pour déterminer s'il doit être présent, on utilise la structure de l'assemblage.

\begin{exemple}
Une colonne atteint une hauteur de $3$ cubes. On ne voit que le cube du sommet et une partie du cube intermédiaire.

La hauteur annoncée impose pourtant :
\[
3
\]
cubes dans cette colonne.

Le cube du bas doit donc être compté, même s'il est entièrement caché.
\end{exemple}

\begin{propriete}
Une représentation en perspective donne des indices visuels, mais le nombre de cubes se déduit de l'organisation spatiale complète, pas seulement de ce qui est visible.
\end{propriete}

\section{8. Vue de dessus}

Une vue de dessus indique les positions occupées lorsque l'on regarde l'objet verticalement.

Elle permet notamment de savoir :
\begin{itemize}
\item combien de colonnes différentes existent ;
\item quelles positions du sol sont occupées ;
\item comment les colonnes sont disposées les unes par rapport aux autres.
\end{itemize}

En revanche, une simple vue de dessus ne donne pas toujours la hauteur de chaque colonne.

Deux assemblages différents peuvent donc avoir exactement la même vue de dessus.

\section{9. Vue de face et vue de côté}

Une vue de face montre les hauteurs observées depuis une direction choisie.

Une vue de côté fournit une information complémentaire.

\coursefigure{/images/mathematiques/6eme/espace-vues-cubes.svg}{Deux représentations planes d'un assemblage : vue de dessus et vue de face.}{Croiser plusieurs vues permet de retrouver davantage d'informations qu'une seule représentation.}

\begin{remarque}
Le mot « face » dans « vue de face » indique une direction d'observation. Il ne désigne pas nécessairement une face géométrique unique du solide.
\end{remarque}

\section{10. Une seule vue peut être insuffisante}

Considérons deux colonnes placées l'une derrière l'autre.

Depuis l'avant, elles peuvent se superposer visuellement.

Une vue de face peut alors montrer une seule colonne de hauteur $3$, alors que plusieurs organisations sont possibles :
\[
3\text{ devant et }1\text{ derrière},
\]
ou
\[
2\text{ devant et }3\text{ derrière},
\]
par exemple.

La vue de dessus permet de savoir combien de positions sont réellement occupées.

\begin{propriete}
Pour reconstituer un assemblage sans ambiguïté, il faut souvent croiser plusieurs vues ou disposer d'informations supplémentaires sur les hauteurs.
\end{propriete}

\section{11. Passer de l'espace au plan}

Passer d'un assemblage réel à une représentation plane consiste à sélectionner une information.

\begin{itemize}
\item La perspective cherche à montrer le relief.
\item La vue de dessus montre l'occupation au sol.
\item La vue de face montre les hauteurs selon une direction.
\item La vue de côté apporte un autre profil de hauteur.
\end{itemize}

Aucune représentation n'est « meilleure » dans tous les cas. On choisit celle qui répond à la question posée.

\section{12. Passer du plan à l'espace}

La démarche inverse demande davantage de raisonnement.

On peut procéder ainsi.

\begin{methode}
Pour reconstituer un assemblage :
\begin{enumerate}
\item utiliser la vue de dessus pour placer les colonnes ;
\item utiliser la vue de face pour contraindre leurs hauteurs ;
\item utiliser la vue de côté pour lever les ambiguïtés ;
\item construire mentalement ou matériellement une solution ;
\item vérifier que toutes les vues correspondent à la construction.
\end{enumerate}
\end{methode}

Une solution qui respecte une vue mais contredit une autre doit être rejetée.

\section{13. Exemple développé : retrouver un nombre de cubes}

Une vue de dessus montre quatre positions occupées. Les informations de hauteur sont :
\[
1,\quad 2,\quad 2,\quad 3.
\]

Le nombre total de cubes est la somme des hauteurs :
\[
N=1+2+2+3.
\]

Ainsi :
\[
\boxed{N=8}.
\]

On peut contrôler le résultat couche par couche.

La première couche contient quatre cubes.

La deuxième couche contient les trois colonnes de hauteur au moins $2$ :
\[
3\ \text{cubes}.
\]

La troisième couche contient uniquement la colonne de hauteur $3$ :
\[
1\ \text{cube}.
\]

Donc :
\[
4+3+1=8.
\]

Les deux stratégies donnent le même résultat.

\section{14. Exemple développé : plusieurs assemblages possibles}

Une vue de dessus montre trois positions alignées.

Une vue de face indique une hauteur maximale de $2$.

Ces données ne suffisent pas forcément à déterminer les trois hauteurs.

Les triplets suivants peuvent par exemple respecter cette hauteur maximale :
\[
(2,1,1)
\]
et
\[
(1,2,1).
\]

Une vue de côté ou une indication supplémentaire peut être nécessaire pour distinguer les assemblages.

\begin{attention}
« Reconstituer un assemblage » ne signifie pas toujours qu'il existe une unique solution. Il faut savoir signaler lorsque les données sont insuffisantes.
\end{attention}

\section{15. Utiliser un matériel réel ou virtuel}

La vision dans l'espace se développe par manipulation.

On peut :
\begin{itemize}
\item empiler de vrais cubes ;
\item tourner l'assemblage ;
\item le dessiner sous plusieurs angles ;
\item comparer le dessin à l'objet ;
\item utiliser un logiciel de géométrie ou de modélisation en trois dimensions.
\end{itemize}

Le passage répété entre l'objet, la représentation et la description verbale est plus important qu'un dessin parfait.

\section{16. Méthode générale}

\begin{methode}
Face à un assemblage de cubes :
\begin{enumerate}
\item repérer le niveau du sol ;
\item identifier les positions occupées ;
\item chercher les hauteurs des colonnes ;
\item choisir un comptage par colonnes ou par couches ;
\item tenir compte des cubes cachés ;
\item confronter plusieurs vues si elles sont disponibles ;
\item vérifier que le total est compatible avec toutes les informations.
\end{enumerate}
\end{methode}

\section{17. Erreurs fréquentes}

\begin{itemize}
\item Compter seulement les cubes visibles.
\item Confondre une face du cube avec une vue du solide.
\item Mesurer directement sur une perspective comme si elle conservait toutes les longueurs.
\item Croire qu'une vue de dessus indique automatiquement les hauteurs.
\item Déduire une solution unique alors que plusieurs assemblages sont compatibles avec les vues.
\item Compter deux fois un même cube en changeant de stratégie.
\end{itemize}

\section{À retenir}

\begin{aretenir}
Le travail essentiel en 6e consiste à \textbf{voir dans l'espace des assemblages de cubes}.

Une perspective donne une image du relief mais peut cacher des cubes.

Une vue de dessus renseigne sur les positions occupées ; une vue de face ou de côté renseigne sur les hauteurs observées.

Pour dénombrer les cubes, on peut additionner les hauteurs des colonnes ou compter les cubes couche par couche.

Une seule vue peut être insuffisante : il faut alors croiser plusieurs représentations et vérifier que la construction proposée satisfait toutes les informations.
\end{aretenir}
